\chapter{The \texorpdfstring{$\aRTS$}{aRTS} family of designs}
\label{chap:design}

This chapter defines the $\alpha$-Rebalancing Targeting Strategies ($\aRTS$),
the central object of the paper, and records the example designs that belong to
the family. It corresponds to Section~3.1 of the paper.

\section{Definition of the family}
\label{sec:design_def}

\begin{definition}[$\aRTS$ family]\label{def:aRTS}
\uses{def:rar, def:plugin_target, def:counts}
Let $K \ge 2$, $\alpha \in [0,1)$ and $m_0 \in \N$ (a burn-in parameter).
A RAR procedure is an \emph{$\alpha$-Rebalancing Targeting Strategy} ($\aRTS$)
if its allocation probabilities $p_{m+1} \in \simplex_K$ satisfy, for all
$m \ge K m_0$, the \emph{throttling condition}
\begin{equation}\label{eq:throttle}
  \forall k \in [K], \quad
  \frac{N_{m,k}}{m} > \hat{\rho}_{m,k}
  \ \Longrightarrow\
  p_{m+1,k} \le \alpha\, \hat{\rho}_{m,k}.
\end{equation}
In words: whenever arm $k$ is currently over-sampled relative to its estimated
target, the probability of selecting it next is throttled down to at most
$\alpha \hat{\rho}_{m,k}$.
\end{definition}

\begin{remark}\label{rem:burnin}
The burn-in parameter $m_0$ is the number of initial samples of each arm that
may be taken before adaptive randomization starts. It plays no role in the
proofs; one may take $m_0 = 0$ when the regularization $\theta_0$ of the
estimator is chosen appropriately.
\end{remark}

\textbf{Lean remark.} The throttling condition \eqref{eq:throttle} is a
\emph{property} of a RAR procedure, not a construction. The cleanest Lean design
is to define a predicate \texttt{IsAlphaRTS} on RAR procedures and prove all
asymptotic theorems from that predicate, then verify it for each concrete
example below. This is realized by \texttt{AlphaRAR.IsARTS} (the throttle predicate) and, for
constructing concrete designs, the shared plumbing
\texttt{AlphaRAR.aRTSAlgorithmOfProb}: it turns any measurable, history-dependent
\emph{probability-vector} function $p$ into an \texttt{Algorithm} whose policy is the finite-support
kernel $\sum_a p_a \delta_a$ (\texttt{AlphaRAR.probVecKernel}); the companion
\texttt{AlphaRAR.aRTSAlgorithmOfProb\_isARTS} reduces verifying \texttt{IsARTS} for a design to the
pure arithmetic of its throttle inequality (\texttt{AlphaRAR.probVecKernel\_apply\_singleton}
identifies the policy mass on arm $k$ with $p_k$). Each concrete example below is then a
probability-vector function whose simplex and throttle properties are elementary.

\section{Examples}
\label{sec:design_examples}

Each of the following designs is shown to belong to the $\aRTS$ family, i.e.\ to
satisfy \eqref{eq:throttle}. These lemmas are independent of one another and are
good small formalization targets.

\begin{definition}[Distance-based design]\label{def:distance_design}
\lean{AlphaRAR.distanceProb, AlphaRAR.distanceAlgorithm}\leanok
\uses{def:plugin_target, def:counts}
Given $\alpha \in [0,1)$, the \emph{distance-based} design sets, for
$m \ge K m_0$ and $k \in [K]$,
\[
  p_{m+1,k} = \alpha\, \hat{\rho}_{m,k}
  + (1-\alpha)\frac{\delta_{m,k}}{\sum_{i=1}^K \delta_{m,i}},
  \qquad
  \delta_{m,k} := \max\!\Big(0,\ \hat{\rho}_{m,k} - \frac{N_{m,k}}{m}\Big),
\]
with the convention that if $\delta_{m,j} = 0$ for all $j$, then
$p_{m+1,j} = \hat{\rho}_{m,j}$ for all $j$.
\end{definition}

\begin{lemma}[Distance-based is a valid probability vector]\label{lem:distance_simplex}
\lean{AlphaRAR.distanceProb_nonneg, AlphaRAR.distanceProb_sum}\leanok
\uses{def:distance_design}
For all $m \ge K m_0$, the vector $p_{m+1}$ of Definition~\ref{def:distance_design}
lies in $\simplex_K$.
\end{lemma}

\begin{proof}\leanok
Each coordinate is nonnegative, and
$\sum_k p_{m+1,k} = \alpha \sum_k \hat{\rho}_{m,k}
 + (1-\alpha)\frac{\sum_k \delta_{m,k}}{\sum_i \delta_{m,i}} = \alpha + (1-\alpha) = 1$
when some $\delta_{m,i}>0$; in the degenerate case $p_{m+1}=\estTarget_m \in \simplex_K$.
\end{proof}

\begin{lemma}[Distance-based is $\aRTS$]\label{lem:distance_isARTS}
\lean{AlphaRAR.distance_isARTS}\leanok
\uses{def:distance_design, def:aRTS}
The distance-based design satisfies the throttling condition \eqref{eq:throttle},
hence belongs to the $\aRTS$ family.
\end{lemma}

\begin{proof}\leanok
\uses{lem:distance_simplex}
Membership in $\simplex_K$ is Lemma~\ref{lem:distance_simplex}. For the throttle:
if $N_{m,k}/m > \hat{\rho}_{m,k}$ then $\delta_{m,k} = 0$, so
$p_{m+1,k} = \alpha\hat{\rho}_{m,k} \le \alpha \hat{\rho}_{m,k}$.

\emph{Formalization.} The design is \texttt{AlphaRAR.distanceProb} (a probability-vector function of
the history) and the algorithm \texttt{AlphaRAR.distanceAlgorithm}; \texttt{AlphaRAR.distanceProb\_%
throttle} is the throttle inequality. The only subtlety is that the mixing normaliser $\sum_i
\delta_{m,i}$ is positive whenever some arm is over-sampled: the signed deficits
$N_{m,i}/m - \hat\rho_{m,i}$ sum to $0$ (\texttt{sum\_histProp}, \texttt{sum\_histTarget}), so an
over-sampled arm forces an under-sampled one, giving $\delta_{m,j}>0$ for some $j$.
\end{proof}

\begin{definition}[ERADE 2025 design]\label{def:erade2025}
\lean{AlphaRAR.eradeProb, AlphaRAR.eradeAlgorithm}\leanok
\uses{def:plugin_target, def:counts}
Given $\alpha \in [0,1)$, the \emph{ERADE 2025} design of
\cite{alkhnefr2025efficient} sets, for $m \ge K m_0$ and $k \in [K]$,
\[
  p_{m+1,k} =
  \begin{cases}
    \alpha \hat{\rho}_{m,k}, & N_{m,k}/m > \hat{\rho}_{m,k},\\
    \hat{\rho}_{m,k}, & N_{m,k}/m = \hat{\rho}_{m,k},\\
    \hat{\rho}_{m,k} + (1-\alpha)\dfrac{\sum_{j\in S}\hat{\rho}_{m,j}}{|T|},
      & N_{m,k}/m < \hat{\rho}_{m,k},
  \end{cases}
\]
where $S = \{ j : N_{m,j}/m > \hat{\rho}_{m,j}\}$ is the set of over-sampled arms
and $T = \{ j : N_{m,j}/m < \hat{\rho}_{m,j}\}$ the set of under-sampled arms.
\end{definition}

\begin{lemma}[ERADE 2025 is a valid probability vector]\label{lem:erade2025_simplex}
\lean{AlphaRAR.eradeProb_nonneg, AlphaRAR.eradeProb_sum}\leanok
\uses{def:erade2025}
For all $m \ge K m_0$, the vector $p_{m+1}$ of Definition~\ref{def:erade2025}
lies in $\simplex_K$.
\end{lemma}

\begin{proof}\leanok
Nonnegativity is clear. For the sum, the total mass removed from over-sampled
arms is $\sum_{j\in S}(1-\alpha)\hat{\rho}_{m,j}$, which is exactly the mass
added on the $|T|$ under-sampled arms, and arms with equality keep their target;
using $\sum_k \hat{\rho}_{m,k} = 1$ gives $\sum_k p_{m+1,k} = 1$.
\end{proof}

\begin{lemma}[ERADE 2025 is $\aRTS$]\label{lem:erade2025_isARTS}
\lean{AlphaRAR.erade_isARTS}\leanok
\uses{def:erade2025, def:aRTS}
The ERADE 2025 design satisfies \eqref{eq:throttle}, hence belongs to the
$\aRTS$ family.
\end{lemma}

\begin{proof}\leanok
\uses{lem:erade2025_simplex}
Membership in $\simplex_K$ is Lemma~\ref{lem:erade2025_simplex}. For the throttle:
if $N_{m,k}/m > \hat{\rho}_{m,k}$ then $p_{m+1,k} = \alpha\hat{\rho}_{m,k}$.

\emph{Formalization.} \texttt{AlphaRAR.eradeProb}/\texttt{eradeAlgorithm} and
\texttt{AlphaRAR.eradeProb\_throttle}. The simplex identity uses that the mass freed from the
over-sampled arms, $\sum_{j\in S}(1-\alpha)\hat\rho_{m,j}$, is redistributed exactly over the $|T|$
under-sampled arms; the degenerate case $|T|=0$ forces $S=\varnothing$ by mass conservation
(\texttt{sum\_histProp}, \texttt{sum\_histTarget}).
\end{proof}

\begin{definition}[Interpolated D-Tracking]\label{def:d_tracking}
\lean{AlphaRAR.dTrackingProb, AlphaRAR.dTrackingAlgorithm}\leanok
\uses{def:plugin_target, def:counts}
For $\alpha \in [0,1)$, the \emph{Interpolated D-Tracking} rule sets, for
$m \ge K m_0$ and $k \in [K]$,
\[
  p_{m+1,k} = \alpha\, \hat{\rho}_{m,k} + (1-\alpha)\,\indic_{\{k = k^\star\}},
  \qquad
  k^\star \in \argmax_{a \in [K]} \big\{ m\hat{\rho}_{m,a} - N_{m,a} \big\}.
\]
The quantity $m\hat{\rho}_{m,a} - N_{m,a}$ is the deficit of arm $a$ relative to
its target, so $k^\star$ is the most under-sampled arm.
\end{definition}

\begin{lemma}[Interpolated D-Tracking is a valid probability vector]\label{lem:d_tracking_simplex}
\lean{AlphaRAR.dTrackingProb_nonneg, AlphaRAR.dTrackingProb_sum}\leanok
\uses{def:d_tracking}
For all $m \ge K m_0$, $p_{m+1} \in \simplex_K$.
\end{lemma}

\begin{proof}\leanok
$\sum_k p_{m+1,k} = \alpha \sum_k \hat{\rho}_{m,k} + (1-\alpha) = 1$, and each
coordinate is nonnegative.
\end{proof}

\begin{lemma}[Interpolated D-Tracking is $\aRTS$]\label{lem:d_tracking_isARTS}
\lean{AlphaRAR.dTracking_isARTS}\leanok
\uses{def:d_tracking, def:aRTS}
For every $\alpha \in [0,1)$, the Interpolated D-Tracking rule satisfies
\eqref{eq:throttle}, hence belongs to the $\aRTS$ family.
\end{lemma}

\begin{proof}\leanok
\uses{lem:counts_sum, lem:d_tracking_simplex}
Membership in $\simplex_K$ is Lemma~\ref{lem:d_tracking_simplex}.
If $N_{m,k}/m > \hat{\rho}_{m,k}$ then $m\hat{\rho}_{m,k} - N_{m,k} < 0$. Since
$\sum_a (m\hat{\rho}_{m,a} - N_{m,a}) = 0$ by Lemma~\ref{lem:counts_sum}, the
maximizing deficit is $\ge 0$, so $k \ne k^\star$ and
$p_{m+1,k} = \alpha\hat{\rho}_{m,k}$.

\emph{Formalization.} The favoured arm $k^\star=\arg\max_a(m\hat\rho_{m,a}-N_{m,a})$ is
\texttt{argmax} of \texttt{AlphaRAR.dtDeficit} (LeanMachineLearning's measurable finite-tuple
\texttt{argmax}, so the policy kernel is measurable without needing an order on the arms); the
throttle is \texttt{AlphaRAR.dTrackingProb\_throttle}.
\end{proof}

\begin{remark}\label{rem:d_tracking_zero}
When $\alpha = 0$, the Interpolated D-Tracking rule becomes
$p_{m+1,k} = \indic_{\{k = k^\star\}}$, which is the D-Tracking rule of
\cite{garivier2016optimal} (without its forced-exploration component, treated in
Chapter~\ref{chap:forced}).
\end{remark}
